\section*{Introducción}

Lo peculiar de esta entrevista es que extiende ante nosotros, casi por completo, la evolución de las preguntas de un investigador de primer nivel a lo largo de más de una década. Desde la clase ACM de la Universidad Jiao Tong de Shanghái, las prácticas de licenciatura y la formación doctoral, hasta FAIR, el aprendizaje autosupervisado, DiT, Cambrian, JEPA y AMI Labs, Xie Saining (谢赛宁) no está narrando una serie de logros aislados entre sí, sino que responde una y otra vez a la misma pregunta: si de verdad queremos entender la inteligencia, ¿en qué variables deberíamos poner la atención? Por eso, esta nota no sigue el orden cronológico, sino que se reorganiza en diez temas, con la intención de condensar en una estructura más clara aquellos juicios de la entrevista que, aunque distribuidos en años distintos, se corresponden entre sí en su lógica.

Si esta entrevista se tomara solo como una «entrevista a un científico estrella», se perdería su parte más valiosa. Lo verdaderamente inusual es que Xie Saining está dispuesto a devolverle a muchas experiencias que podrían presentarse de forma más pulida su condición original: un proceso lleno de azar, fracasos, errores de juicio, retrasos y correcciones repetidas. Cuando habla de su asesor, de sus artículos, de emprender y de los modelos de mundo, desmonta una y otra vez el mismo malentendido: lo que ve el público es la superficie lisa del resultado, mientras que a lo que el investigador realmente tiene que enfrentarse es a ese caos sumamente irregular que precede al resultado. Por eso insiste tanto en la no linealidad, la exploración, los resultados negativos y la antifragilidad, pues estos conceptos son precisamente las condiciones necesarias para que de ese caos surja una estructura nueva.

El segundo nivel de valor de la entrevista está en que le devuelve al «estudio de la visión» su peso intelectual. En el contexto de los modelos grandes, muchos entienden, sin darse cuenta, vision como un componente periférico del modelo de lenguaje, o entienden world model como un nuevo empaque narrativo. Pero el argumento de Xie Saining es muy claro: la visión importa no porque pueda agregarle al LLM una modalidad más, sino porque conecta con el mundo real, con el espacio continuo, con el flujo del tiempo y con las consecuencias de las acciones. Dicho de otro modo, si se entiende la inteligencia únicamente como generar el siguiente símbolo en el espacio de tokens, entonces todo el campo de la IA confunde con facilidad «saber decir» con «saber entender».

\begin{importantbox}{Tres hilos conductores para leer esta nota}
El primer hilo conductor es «el proceso de convertirse en investigador», que incluye cómo la familia, la escuela, el asesor y los pares moldean juntos a la persona. El segundo hilo conductor es «la continuidad teórica que va de la visión al modelo de mundo», que explica por qué Xie Saining nunca ha visto su propio trabajo como una serie de saltos discretos. El tercer hilo conductor es «cómo se sostienen mutuamente la investigación y la organización»: de FAIR a AMI Labs, lo que en realidad cambia no es el problema en sí, sino el entorno institucional que lo sostiene.
\end{importantbox}

Por último, esta nota conserva deliberadamente la parte cultural y vital de la entrevista. El cine, la filosofía, Nueva York, la inteligencia animal, Wittgenstein, Klopp, las listas de libros y el sentido del destino no son adornos de la discusión técnica, sino un trasfondo inseparable de cómo Xie Saining entiende la inteligencia. Sin ese trasfondo, su research taste, su cautela ante los límites del lenguaje y su insistencia en el mundo real se malinterpretarían como unas cuantas frases vistosas. Precisamente porque él siempre mide la técnica contra un mundo vital más amplio, esta entrevista no habla solo de IA, sino también de cómo mantenerse lúcido en una época de competencia acelerada y narrativas sobrecalentadas.

